\documentclass[11pt]{article}
\usepackage[a4paper,margin=1.9cm]{geometry}
\usepackage[T1]{fontenc}
\usepackage{textcomp}
\usepackage[spanish]{babel}
\usepackage{amsmath,amssymb}
\usepackage{enumitem}
\usepackage{tikz}
\usetikzlibrary{calc,arrows.meta,patterns,decorations.pathmorphing}
\usepackage{xcolor}
\usepackage{multicol}
\usepackage{parskip}
\usepackage{lmodern}
\renewcommand{\familydefault}{\sfdefault}

\definecolor{unalblue}{RGB}{31,73,124}

\newlist{opciones}{enumerate}{1}
\setlist[opciones]{label=(\alph*),itemsep=1pt,topsep=2pt,leftmargin=1.8em,font=\normalfont}

\newcommand{\vect}[2]{\begin{pmatrix}#1\\#2\end{pmatrix}}
\newcommand{\vectiii}[3]{\begin{pmatrix}#1\\#2\\#3\end{pmatrix}}
\newcommand{\vectiv}[4]{\begin{pmatrix}#1\\#2\\#3\\#4\end{pmatrix}}

\newcommand{\examheader}{%
  \noindent
  \begin{minipage}[t]{0.6\linewidth}
    {\small\bfseries Álgebra Lineal --- Simulacro Primer Examen Parcial}\\
  \end{minipage}%
  \hfill
  \begin{minipage}[t]{0.35\linewidth}
    \raggedleft\small Escuela de Matemáticas. Universidad Nacional de Colombia, Sede Medellín.
  \end{minipage}\\[2pt]
  \rule{\linewidth}{0.6pt}
}
\newcommand{\examfooter}{%
  \vfill
  \rule{\linewidth}{0.4pt}\\[1pt]
  {\scriptsize\textit{Transcripción digital --- MathUNAL}\hfill\textcolor{unalblue}{\textbf{\thepage}}}
}
\pagestyle{empty}

\begin{document}

\examheader

\begin{center}
{\bfseries\large Simulacro Primer Examen Parcial de Álgebra Lineal}
\end{center}

\textbf{Problema 1 (10 puntos):} Considere el sistema de ecuaciones lineales en las variables $x,y,z,w$
\begin{align*}
x + 2y - 3z &= 0\\
y + 5w &= 0\\
2x + z - w &= 0\\
y + kw &= 0
\end{align*}
Encontrar todos los valores de $k$ tales que el sistema tiene infinitas soluciones.

\vspace{2cm}

\examfooter
\newpage

\examheader

\textbf{Problema 2 (10 puntos):} Sean $v$ y $w$ dos vectores de $\mathbb{R}^3$ con
\[
v = \vectiii{-5}{-1}{-5} \quad \text{y} \quad w = \vectiii{18}{-30}{-12}.
\]
Determine el valor de $\lambda > 0$ tal que
\[
\|\lambda v + w\|^2 = 2184.
\]

\vspace{2cm}

\examfooter
\newpage

\examheader

\textbf{Problema 3 (10 puntos):} Sean $v = \vectiii{-8}{17}{0}$ y $w = \vectiii{1}{\lambda}{11}$ dos vectores en $\mathbb{R}^3$. Encontrar el valor de $\lambda$ para que el vector $x = \vectiii{53}{-32}{55}$ pertenezca al plano generado por los vectores $v$ y $w$.

\vspace{2cm}

\examfooter
\newpage

\examheader

\textbf{Problema 4 (10 puntos):} Dados dos vectores columna $u = \vectiv{-5}{1}{-7}{6}$ y $v = \vectiv{-7}{8}{6}{6}$ calcule $uv^T$ y el rango de $uv^T$.

\vspace{2cm}

\examfooter
\newpage

\examheader

\textbf{Problema 5 (10 puntos):} Determinar si los siguientes vectores generan a $\mathbb{R}^3$
\[
\vectiii{1}{1}{1}, \quad \vectiii{1}{1}{0}, \quad \vectiii{0}{0}{-3}, \quad \vectiii{1}{1}{5}.
\]

\vspace{2cm}

\examfooter
\newpage

\examheader

\textbf{Problema 6 (10 puntos):} Supongamos que
\[
A = \begin{pmatrix} a_{11} & a_{12} & a_{13} \\ a_{21} & a_{22} & a_{23} \\ a_{31} & a_{32} & a_{33} \end{pmatrix}
\]
es una matriz de tamaño $3\times 3$ tal que $\det(A) = -7$. Sean $B$ y $C$ las matrices dadas por
\[
B = \begin{pmatrix} a_{11} & 9a_{12} & a_{13} \\ 10a_{21} & 90a_{22} & 10a_{23} \\ a_{31} & 9a_{32} & a_{33} \end{pmatrix}, \qquad
C = \begin{pmatrix} 1 & 0 & -7 \\ 2 & 0 & 1 \\ 8 & 5 & -3 \end{pmatrix}.
\]
Si $D = B^T C$, calcular el determinante de $D$.

\vspace{2cm}

\examfooter
\newpage

\examheader

\textbf{Problema 7 (10 puntos):} Dados los vectores
\[
v_4 = \vectiv{-7}{9}{6}{0}, \quad v_1 = \vectiv{4}{0}{3}{8}, \quad v_5 = \vectiv{21}{18}{-9}{10}, \quad v_2 = \vectiv{6}{9}{0}{6}, \quad v_3 = \vectiv{0}{0}{0}{0},
\]
señale todos los conjuntos que son Linealmente Independientes. Puede haber más de una respuesta.

\begin{opciones}
\item $\{v_1, v_2\}$
\item $\{v_1, v_2, v_3\}$
\item $\{v_1, v_2, v_4\}$
\item $\{v_1, v_2, v_4, v_5\}$
\item $\{v_1, v_2, v_3, v_4\}$
\item $\{v_2, v_4\}$
\end{opciones}

\examfooter

\end{document}
